% =============================================================================
% Section 5: Baseline Framework Evaluation
% =============================================================================

\paragraph{Scope note.} This section gives compact framework-level context only. It is separate from the primary cross-family grounding-effect study in \S\ref{sec:gisr}; detailed framework, BIRD, DIN-SQL, token-cost, and case-study tables are retained in Supplement~S1. The $+32.50$ pp grounding effect on Robustness reported in \S\ref{sec:gisr} derives from the 40-question CQ-125 subset and is \emph{not} the same experiment as the earlier 15-question pilot noted below.

\subsection{Supporting framework results}
\label{sec:gis-results}

In the earlier agent-loop GIS run, the best observed framework sample reaches Spatial EX $0.706$ on the 85-question set with a \texttt{gemini-2.5-flash} backbone, versus a schema-only baseline of $0.518$. This best-sample result is descriptive and is not used as primary inferential evidence. On the 15-question Robustness pilot, the framework reaches $7/7$ on discordant pairs ($p=0.0156$ exact). On BIRD held-out 150q, the Full pipeline reaches EX $0.507$ vs.\ DIN-SQL $0.440$ ($p=0.0755$, marginal). On the 40-question Robustness suite, the gap concentrates on safety/refusal categories where DIN-SQL has no built-in handling (Full $0.975$ vs.\ DIN-SQL $0.275$). These supporting results motivate the framework context; the main CQ-125 claims use the question-level majority-vote protocol in \S\ref{sec:gisr}.

\subsection{Gemma4 model-size sweep on a fixed local host}
\label{sec:gemma4-scale-sweep}

As a deployment-oriented probe, we ran five Gemma4 model scales on one fixed local Ollama serving node using CQ-125 and the same baseline-vs-full framework configurations. The source table is provided as \texttt{gemma4\_host228\_scale\_sweep\_summary.csv}. The serving metadata were queried from the same node via the Ollama \texttt{/api/tags} and \texttt{/api/version} endpoints: Ollama version 0.30.5; GGUF format; Q4\_K\_M quantisation for all five Gemma4 models; reported parameter sizes of 5.1B, 8.0B, 11.9B, 25.8B, and 31.3B. The semantic retrieval backend on the same node used \texttt{nomic-embed-text-v2-moe}. The full model-metadata source table is distributed as \texttt{host228\_ollama\_model\_metadata.csv}. This probe is descriptive: it is not an inferential cross-family test and is not used in the primary majority-vote McNemar reductions.

In this fixed-host summary, full-pipeline EX ranged from 80/125 to 114/125, compared with 37/125 to 71/125 for schema-only generation. The \texttt{Gemma4:31b} full run had the highest observed EX (114/125, 20.04 min), while \texttt{Gemma4:26b} was one question lower (113/125, 12.50 min). These numbers are reported only to illustrate an accuracy--runtime deployment trade-off. Because several host228 full-pipeline runs existed for some scales, the CSV identifies the selected full rows as best-observed engineering summaries rather than stochastic significance evidence.

\subsection{Reproducibility}
\label{sec:reproducibility-stub}

All commands and artefacts to regenerate the previously released framework-level numbers are in the anonymised reproducibility repository (\url{https://anonymous.4open.science/r/nl2geosql-reproduction-36D6/}, $\approx$76 JSONL plus offline analysis scripts; reviewer runtime $\approx 90$\,s with zero LLM API and zero PostgreSQL dependency). The cross-family analysis in \S\ref{sec:gisr} is regenerated from the same repository. The newer Gemma4 deployment probe is reported from frozen local run summaries; its source-table summary is included with this manuscript source as \texttt{gemma4\_host228\_scale\_sweep\_summary.csv} and is not used in the primary statistical reductions.
